\chapter{Regulasi, Etika, \& Kebijakan AI Global}
\label{chap:regulation}

Seiring dengan kemajuan pesat teknologi AI, pemerintah di seluruh dunia berlomba-lomba untuk membuat kerangka hukum yang mengatur penggunaannya. Appendix ini merangkum regulasi utama yang perlu diketahui oleh praktisi AI.

\section{Uni Eropa: The EU AI Act}
Regulasi AI paling komprehensif di dunia saat ini, yang mengadopsi pendekatan berbasis risiko (risk-based approach).

\subsection{Kategori Risiko}
\begin{enumerate}
    \item \textbf{Unacceptable Risk (Dilarang):}
    \begin{itemize}
        \item Sistem skor sosial (social scoring) oleh pemerintah.
        \item Identifikasi biometrik real-time di ruang publik oleh penegak hukum (dengan pengecualian ketat).
        \item Manipulasi perilaku subliminal yang membahayakan.
        \item Eksploitasi kerentanan kelompok tertentu (anak-anak, penyandang disabilitas).
    \end{itemize}
    \item \textbf{High Risk (Diatur Ketat):}
    \begin{itemize}
        \item Infrastruktur kritis (transportasi, energi).
        \item Pendidikan dan pelatihan vokasi (menentukan akses/penerimaan).
        \item Ketenagakerjaan (penyaringan CV otomatis).
        \item Layanan publik dan swasta esensial (kredit skor, asuransi).
        \item Penegakan hukum dan imigrasi.
    \end{itemize}
    \textbf{Kewajiban untuk High Risk AI:}
    \begin{itemize}
        \item Penilaian kesesuaian (conformity assessment) sebelum masuk pasar.
        \item Tata kelola data yang ketat (mencegah bias).
        \item Dokumentasi teknis dan pencatatan (logging) otomatis.
        \item Transparansi bagi pengguna.
        \item Pengawasan manusia (human oversight).
        \item Keamanan siber dan akurasi tinggi.
    \end{itemize}
    \item \textbf{Limited Risk (Transparansi Khusus):}
    \begin{itemize}
        \item Chatbot (harus memberitahu user bahwa mereka bicara dengan mesin).
        \item Deepfakes (harus diberi label bahwa konten tersebut dimanipulasi).
    \end{itemize}
    \item \textbf{Minimal Risk (Bebas):}
    \begin{itemize}
        \item Sebagian besar aplikasi AI saat ini (filter spam, video game).
    \end{itemize}
\end{enumerate}

\section{Amerika Serikat: Executive Order on AI}
Pada Oktober 2023, Presiden Joe Biden menandatangani Perintah Eksekutif tentang Pengembangan dan Penggunaan AI yang Aman, Terjamin, dan Terpercaya.

\subsection{Poin Kunci}
\begin{itemize}
    \item \textbf{Standar Keamanan:} Mewajibkan pengembang model AI "dual-use" yang sangat besar (seperti GPT-5 atau lebih besar) untuk membagikan hasil uji keamanan (red-teaming) kepada pemerintah AS sebelum dirilis.
    \item \textbf{Privasi:} Memperkuat perlindungan privasi dengan mendanai riset teknologi pelestarian privasi (privacy-preserving technologies).
    \item \textbf{Hak Sipil:} Mencegah diskriminasi algoritma dalam perumahan, pekerjaan, dan kredit.
    \item \textbf{Watermarking:} Mengembangkan standar untuk menandai konten yang dihasilkan AI agar masyarakat bisa membedakannya dari konten asli.
\end{itemize}

\section{Tiongkok: Generative AI Measures}
Tiongkok telah mengeluarkan serangkaian regulasi yang menargetkan algoritma rekomendasi, deep synthesis (deepfakes), dan generative AI.

\subsection{Poin Kunci}
\begin{itemize}
    \item \textbf{Nilai Sosialis:} Konten yang dihasilkan AI harus sejalan dengan nilai-nilai inti sosialis.
    \item \textbf{Pendaftaran Algoritma:} Penyedia layanan harus mendaftarkan algoritma mereka ke pemerintah.
    \item \textbf{Labeling:} Kewajiban ketat untuk memberi label pada konten sintetis.
    \item \textbf{Tanggung Jawab:} Penyedia layanan bertanggung jawab atas konten yang dihasilkan oleh platform mereka.
\end{itemize}

\section{Indonesia: Surat Edaran Menkominfo No. 9 Tahun 2023}
Indonesia mengambil pendekatan "soft law" melalui Surat Edaran tentang Etika Kecerdasan Artifisial.

\subsection{Nilai Etika AI}
\begin{enumerate}
    \item \textbf{Inklusivitas:} AI harus bermanfaat bagi semua kalangan.
    \item \textbf{Kemanusiaan:} Melindungi hak asasi manusia.
    \item \textbf{Keamanan:} Menjamin keamanan data dan sistem.
    \item \textbf{Demokrasi:} Tidak memanipulasi proses demokrasi.
    \item \textbf{Transparansi:} Keterbukaan tentang penggunaan AI.
    \item \textbf{Kredibilitas:} Informasi yang dihasilkan dapat dipercaya.
    \item \textbf{Akuntabilitas:} Ada pihak yang bertanggung jawab jika terjadi kesalahan.
\end{enumerate}
Meskipun tidak memiliki sanksi pidana langsung (karena hanya Surat Edaran), ini menjadi pedoman bagi industri dan bisa menjadi dasar bagi regulasi yang lebih mengikat di masa depan (UU PDP tetap berlaku untuk aspek data pribadi).

\section{Isu Hak Cipta (Copyright) & AI}

\subsection{Input: Apakah melatih AI dengan data berhak cipta itu legal?}
\begin{itemize}
    \item \textbf{Fair Use (AS):} Perusahaan AI (OpenAI, Google) berargumen bahwa pelatihan model adalah "transformative use" yang dilindungi doktrin Fair Use.
    \item \textbf{Gugatan Hukum:} Banyak seniman dan penulis (New York Times, Getty Images) menggugat perusahaan AI karena menggunakan karya mereka tanpa izin atau kompensasi. Hasil kasus-kasus ini akan menentukan masa depan industri.
    \item \textbf{Opt-Out:} Beberapa platform (seperti ArtStation) mulai menawarkan opsi "No AI" agar karya tidak discrape.
\end{itemize}

\subsection{Output: Siapa pemilik hak cipta karya AI?}
\begin{itemize}
    \item \textbf{AS (USCO):} Karya yang murni dihasilkan oleh AI (tanpa input kreatif manusia yang signifikan) \textbf{tidak dapat memiliki hak cipta}. Namun, jika manusia melakukan modifikasi berat, bagian modifikasinya bisa dilindungi.
    \item \textbf{Inggris:} Memberikan perlindungan hak cipta untuk karya "computer-generated", dengan penulis dianggap sebagai orang yang melakukan pengaturan yang diperlukan untuk penciptaan karya tersebut.
\end{itemize}

\section{GDPR & AI (Eropa)}
General Data Protection Regulation (GDPR) memiliki dampak besar pada AI, terutama terkait data pribadi.

\begin{itemize}
    \item \textbf{Right to Explanation (Pasal 22):} Individu berhak untuk tidak tunduk pada keputusan yang semata-mata didasarkan pada pemrosesan otomatis (termasuk profiling) yang memiliki efek hukum. Mereka berhak meminta penjelasan logis di balik keputusan tersebut.
    \item \textbf{Data Minimization:} Hanya mengumpulkan data yang benar-benar diperlukan. Ini bertentangan dengan praktik Deep Learning yang haus data.
    \item \textbf{Right to be Forgotten:} Pengguna berhak meminta data mereka dihapus. Ini sulit dilakukan pada model AI yang sudah dilatih (apakah model harus di-retrain dari nol? Atau bisa dilakukan "machine unlearning"?).
\end{itemize}

\section{Checklist Kepatuhan Etika untuk Developer}

Sebelum meluncurkan produk AI, tanyakan hal-hal berikut:
\begin{itemize}
    \item [ ] **Bias:** Apakah dataset pelatihan mewakili demografi pengguna secara adil?
    \item [ ] **Privasi:** Apakah ada PII (Personally Identifiable Information) yang bocor dalam training data?
    \item [ ] **Transparansi:** Apakah pengguna tahu bahwa mereka berinteraksi dengan AI?
    \item [ ] **Explainability:** Bisakah kita menjelaskan mengapa model membuat keputusan tertentu?
    \item [ ] **Fallback:** Apa yang terjadi jika model gagal atau berhalusinasi? Apakah ada manusia in-the-loop?
\end{itemize}
